% ===================== PREAMBULE COMMUN (à copier VERBATIM en tête de CHAQUE .tex) =====================
\documentclass[11pt,a4paper]{article}
\usepackage[utf8]{inputenc}
\usepackage[T1]{fontenc}
\usepackage{lmodern}
\usepackage[french]{babel}
\usepackage{amsmath,amssymb,mathtools}
\usepackage{icomma}
\usepackage{xcolor}
\usepackage{colortbl}
\usepackage{tikz}
\usepackage{pgfplots}
\usepackage[most]{tcolorbox}
\usepackage{enumitem}
\usepackage{array}
\usepackage{booktabs}
\usepackage{multicol}
\usepackage{fancyhdr}
\usepackage{caption}
\captionsetup{labelformat=empty,justification=centering,font=small}
\usepackage[hidelinks]{hyperref}
\usetikzlibrary{arrows.meta,positioning,calc,angles,quotes,patterns,backgrounds,shapes.geometric}
\pgfplotsset{compat=1.18}
\usepackage[a4paper,margin=2.2cm,headheight=26pt,headsep=10pt]{geometry}

\definecolor{primary}{RGB}{219,54,77}
\definecolor{primarydark}{RGB}{150,28,46}
\definecolor{primarylight}{RGB}{250,224,228}
\definecolor{defcol}{RGB}{0,114,178}
\definecolor{thmcol}{RGB}{0,140,120}
\definecolor{excol}{RGB}{0,135,90}
\definecolor{methcol}{RGB}{90,90,100}
\definecolor{courbe}{RGB}{40,90,200}
\definecolor{courbeB}{RGB}{210,60,40}
\definecolor{courbeC}{RGB}{30,150,80}
\definecolor{axiscol}{RGB}{60,60,60}
\definecolor{gridcol}{RGB}{200,205,215}

\tcbset{boxbase/.style={enhanced,breakable,boxrule=0.9pt,arc=2.5pt,
   left=3mm,right=3mm,top=2.3mm,bottom=2.3mm,
   fonttitle=\bfseries,coltitle=white,toptitle=0.6mm,bottomtitle=0.6mm}}
\newtcolorbox{methode}[1][]{boxbase,colback=methcol!7,colframe=methcol,sharp corners,
  title={\textsc{Méthode}\ifx\relax#1\relax\relax\else\textmd{ --- \itshape #1}\fi}}
\newtcolorbox{exemple}[1][]{boxbase,colback=excol!5,colframe=excol,
  title={\textsc{Exemple}\ifx\relax#1\relax\relax\else\textmd{ : \itshape #1}\fi}}
\newtcolorbox{idee}[1][]{boxbase,colback=defcol!6,colframe=defcol,
  title={\textsc{L'essentiel}\ifx\relax#1\relax\relax\else\textmd{ : \itshape #1}\fi}}
\newtcolorbox{piege}[1][]{boxbase,colback=primary!7,colframe=primary,
  title={\textsc{Piège}\ifx\relax#1\relax\relax\else\textmd{ : \itshape #1}\fi}}

\pgfplotsset{repere/.style={axis lines=middle,
    axis line style={-{Stealth[length=2.2mm]},axiscol,line width=0.8pt},
    label style={font=\footnotesize},tick label style={font=\scriptsize},
    every axis x label/.style={at={(current axis.right of origin)},anchor=north west},
    every axis y label/.style={at={(current axis.above origin)},anchor=south east},
    xlabel={$x$},ylabel={$y$},grid=both,minor tick num=0,
    major grid style={gridcol,line width=0.4pt},samples=200,clip=false}}

\pagestyle{fancy}\fancyhf{}
\fancyhead[L]{\small\textcolor{primarydark}{\textbf{Fiche 4} — Les fonctions}}
\fancyhead[R]{\small\textcolor{primarydark}{\textsc{Carnet d'entraînement}}}
\fancyfoot[C]{\small\thepage}
\renewcommand{\headrulewidth}{0.8pt}
\renewcommand{\headrule}{\hbox to\headwidth{\color{primary}\leaders\hrule height \headrulewidth\hfill}}
\usepackage{titlesec}
\titleformat{\section}{\Large\bfseries\color{primarydark}}{\thesection}{0.6em}{}
\titleformat{\subsection}{\large\bfseries\color{primary}}{\thesubsection}{0.6em}{}

\newcommand{\R}{\mathbb{R}}\newcommand{\Z}{\mathbb{Z}}\newcommand{\N}{\mathbb{N}}
\newcommand{\Df}{D_f}
\newcommand{\croit}{\textcolor{thmcol}{\Large$\nearrow$}}
\newcommand{\decroit}{\textcolor{thmcol}{\Large$\searrow$}}

\newcommand{\exercice}[1]{\medskip\noindent{\color{primarydark}\bfseries\large Exercice #1}\par\nopagebreak\smallskip}
\newcommand{\titrecarnet}[3]{%
\begin{center}\begin{tikzpicture}
\node[rectangle,rounded corners=7pt,fill=primary,text=white,minimum width=15cm,
  minimum height=1.7cm,align=center,inner sep=8pt]
  {{\LARGE\bfseries #1}\\[2pt]{\normalsize #2}};
\end{tikzpicture}\end{center}
\begin{center}\itshape\color{primarydark}#3\end{center}\medskip}

% ---- RÈGLES D'USAGE DES MACROS ----
% * \croit et \decroit s'emploient en mode TEXTE (dans une cellule de tableau),
%   JAMAIS entre $...$ (ils contiennent déjà un $\nearrow$).
% * \titrecarnet{Titre}{Sous-titre}{Ligne en italique} : bandeau de titre.
% * \exercice{1} : titre d'exercice.
% * Boîtes : \begin{idee}...\end{idee}, \begin{exemple}[titre]...\end{exemple},
%   \begin{methode}[titre]...\end{methode}, \begin{piege}[titre]...\end{piege}.
% Le corps du document commence après \begin{document}.
\begin{document}
\titrecarnet{Thème 4 — Fonction du second degré}{Corrigé — niveau Difficile}{Détail des calculs.}

\exercice{1}
\begin{enumerate}[label=\alph*),leftmargin=1.8em,topsep=2pt]
  \item Racines $-1$ et $2$ : $P(x)=a(x-(-1))(x-2)=\boxed{a(x+1)(x-2)}$.
  \item $P(1)=a(1+1)(1-2)=a\times2\times(-1)=-2a$. Or $P(1)=-4$, donc $-2a=-4\Rightarrow \boxed{a=2}$.
  \item $P(0)=a(0+1)(0-2)=2\times1\times(-2)=\boxed{-4}$.
\end{enumerate}

\exercice{2}
On résout $x^2-3x+1=x+1$ :
\[ x^2-3x+1-x-1=0\Leftrightarrow x^2-4x=0\Leftrightarrow x(x-4)=0\Leftrightarrow x=0 \text{ ou } x=4. \]
Ordonnées (via $y=x+1$) : pour $x=0$, $y=1$ ; pour $x=4$, $y=5$. Points d'intersection :
\[ \boxed{(0;1) \text{ et } (4;5)}. \]

\exercice{3}
\begin{enumerate}[label=\alph*),leftmargin=1.8em,topsep=2pt]
  \item Deux largeurs $x$ et une longueur $L$ utilisent le grillage : $2x+L=20\Rightarrow L=20-2x$. Aire :
        \[ A(x)=x\,L=x(20-2x)=-2x^2+20x, \qquad 0<x<10. \]
  \item $A$ est une parabole avec $a=-2<0$ (sommet = maximum). Abscisse du sommet : $x=-\dfrac{b}{2a}=-\dfrac{20}{2\times(-2)}=\boxed{5}$~m.
  \item $A(5)=-2\times25+20\times5=-50+100=\boxed{50\ \text{m}^2}$.
\end{enumerate}

\exercice{4}
$f(x)=x^2-4x+3$, $a=1,\ b=-4,\ c=3$.
\begin{enumerate}[label=\alph*),leftmargin=1.8em,topsep=2pt]
  \item $\Df=\R$ ; $\Delta=16-12=4>0$ ; racines $x=\dfrac{4\pm2}{2}$, soit $x_1=1$ et $x_2=3$.
  \item Sommet $\Big(-\dfrac{b}{2a};-\dfrac{\Delta}{4a}\Big)=(2;-1)$ ; axe $x=2$ ; $Oy:(0;3)$ ; forme factorisée $f(x)=(x-1)(x-3)$.
  \item $a>0$ : décroissante puis croissante, minimum en $x=2$.
\end{enumerate}

\begin{center}
\renewcommand{\arraystretch}{1.6}\setlength{\tabcolsep}{12pt}
\begin{tabular}{|c|ccccc|}
\hline\rowcolor{primarylight}
$x$ & $-\infty$ & & $2$ & & $+\infty$ \\\hline
$f'(x)=2x-4$ & & $-$ & $0$ & $+$ & \\\hline
$f(x)$ & & \decroit & $-1$ & \croit & \\\hline
\end{tabular}
\end{center}

\begin{center}
\begin{tikzpicture}
\begin{axis}[repere,width=9.5cm,height=5.6cm,xmin=-1,xmax=5,ymin=-2,ymax=4,
   xtick={1,2,3,4},ytick={-1,1,2,3},samples=100,title={\footnotesize $f(x)=x^2-4x+3$}]
  \addplot[courbe,line width=1.3pt,domain=-0.45:4.45]{x^2-4*x+3};
  \filldraw[primary](axis cs:2,-1)circle(2pt)node[anchor=north,font=\scriptsize]{sommet $(2;-1)$};
  \filldraw[courbeC](axis cs:1,0)circle(1.6pt)node[anchor=south east,font=\scriptsize]{$1$};
  \filldraw[courbeC](axis cs:3,0)circle(1.6pt)node[anchor=south west,font=\scriptsize]{$3$};
\end{axis}
\end{tikzpicture}
\end{center}
\end{document}
